\documentclass[twocolumn]{aastex631}

\usepackage{amsmath}
\usepackage{multirow}
\usepackage{natbib}
\usepackage{graphicx} 
\usepackage{aas_macros}

\begin{document}

Degenerate Higher-Order Scalar-Tensor (DHOST) theories are constructed to be free of pathological ghost instabilities at the classical level, but their viability under quantum corrections remains an open question. To probe the quantum fate of this classical stability, we isolate and analyze a specific subclass of Type Ia DHOST theories that possesses a field-dependent gauge symmetry, providing a controlled theoretical framework. After deriving the mathematical conditions required for this symmetry to exist, we select a simple representative model and compute its one-loop effective action using the background field method. We demonstrate that the classical symmetry is broken by quantum corrections, resulting in a quantum anomaly. Crucially, the divergent counter-terms that constitute this anomaly, including a Weyl-squared invariant, have a structure that explicitly violates the theory's original degeneracy conditions. This violation reintroduces the very Ostrogradsky ghost that was meticulously eliminated at the classical level, rendering the theory quantum-mechanically unstable. This result establishes a direct link between anomaly generation and quantum instability, constituting a no-go theorem for this class of symmetric theories and elevating anomaly cancellation to a powerful selection principle for constructing consistent theories of modified gravity.

\end{document}
                